\begin{exercise}{8}
Calculate the limit and colimit of a simplicial set
$X\colon\Delta^{\mathrm{op}}\to\mathrm{Set}$.

\begin{proof}
There are canonical identifications
\[
 \lim_{\Delta^{\mathrm{op}}}X\cong X_0,
 \qquad
 \operatorname*{colim}_{\Delta^{\mathrm{op}}}X\cong\pi_0^\Delta(X).
\]
For the limit, use the compatible-family construction of
Exercise~\ref{ex:6}. Let $p_n\colon[n]\to[0]$ be the unique map.
A compatible family $(x_n)$ must satisfy $x_n=p_n^*x_0$.
Conversely, these formulas define a compatible family for every
$x_0\in X_0$, since $p_n\alpha=p_m$ for every
$\alpha\colon[m]\to[n]$. Projection to $x_0$ is therefore a
bijection, with limit cone given by the maps $p_n^*\colon X_0\to X_n$.

By the colimit construction of Exercise~\ref{ex:6}, the colimit is the
quotient of $\coprod_{n\ge0}X_n$ by the identifications
$x\sim\alpha^*x$. Every simplex is thus identified with each of its
vertices. Any two vertices of a simplex belong to the same component
of $\pi_0^\Delta(X)$: restriction along the increasing map
$[1]\to[n]$ selecting them gives an edge between them.
Consequently, sending a simplex to the component of any of its
vertices is well-defined and respects every relation
$x\sim\alpha^*x$. This gives a map from the colimit to
$\pi_0^\Delta(X)$.
Conversely, the map from $X_0$ to the colimit identifies the endpoints
of every edge, so it factors through $\pi_0^\Delta(X)$ by
Definition~1.1.9. The two maps are inverse, since every simplex is
identified with a vertex.
\end{proof}
\end{exercise}
